\documentclass[12pt, letterpaper]{article}

% ── Packages ──────────────────────────────────────────────────────────────────
\usepackage[margin=1in]{geometry}
\usepackage{setspace}
\usepackage{times}           % Times New Roman font
\usepackage{titlesec}        % Custom section headings
\usepackage{parskip}         % Paragraph spacing
\usepackage{hyperref}        % Clickable references
\usepackage{natbib}          % Bibliography (swap for biblatex if preferred)
\usepackage{lipsum}          % REMOVE when filling in real content

% ── Spacing ───────────────────────────────────────────────────────────────────
\doublespacing

% ── Section Formatting ────────────────────────────────────────────────────────
\titleformat{\section}{\large\bfseries}{\Roman{section}.}{1em}{}
\titleformat{\subsection}{\normalsize\bfseries}{\Alph{subsection}.}{1em}{}
\titleformat{\subsubsection}{\normalsize\itshape}{\arabic{subsubsection}.}{1em}{}

% ── Hyperlink Setup ───────────────────────────────────────────────────────────
\hypersetup{
    colorlinks=true,
    linkcolor=black,
    citecolor=black,
    urlcolor=blue
}

% ══════════════════════════════════════════════════════════════════════════════
\begin{document}

% ── Title Page ────────────────────────────────────────────────────────────────
\begin{titlepage}
    \centering
    \vspace*{2cm}
    {\LARGE\bfseries Zero Trust Architecture: Redefining Security\\
    in the Modern Enterprise\par}
    \vspace{2cm}
    {\large Your Name\par}
    \vspace{0.5cm}
    {\large Course Name / Number\par}
    {\large Institution Name\par}
    \vspace{0.5cm}
    {\large \today\par}
    \vfill
\end{titlepage}

% ── Abstract ──────────────────────────────────────────────────────────────────
\begin{abstract}
% ~150 words — brief overview of ZTA, why perimeter-based security is failing,
% what this paper examines, and the key conclusion.
%
% TODO: Write abstract here.
\lipsum[1] % DELETE this line and replace with your abstract
\end{abstract}

\newpage
\tableofcontents
\newpage

% ══════════════════════════════════════════════════════════════════════════════
% I. INTRODUCTION  (~0.5 page)
% ══════════════════════════════════════════════════════════════════════════════
\section{Introduction}

\subsection{Hook / Opening}
% Rise in cyberattacks, data breaches, ransomware — set the scene.

\subsection{Problem Statement}
% The "castle-and-moat" perimeter model no longer works in a cloud-first,
% remote-work world.

\subsection{Thesis Statement}
% Zero Trust Architecture offers a more resilient, identity-centric approach
% to cybersecurity for modern distributed environments.

\subsection{Roadmap}
% Brief overview of what each section covers.

% ══════════════════════════════════════════════════════════════════════════════
% II. BACKGROUND \& HISTORY  (~1 page)
% ══════════════════════════════════════════════════════════════════════════════
\section{Background and History}

\subsection{Traditional Perimeter-Based Security Models}

    \subsubsection{How They Worked}
    % Trusted inside, untrusted outside — the implicit trust model.

    \subsubsection{Limitations}
    % Failure modes exposed by cloud computing, BYOD, and remote work.

\subsection{Origins of Zero Trust}

    \subsubsection{John Kindervag's 2010 Forrester Research Paper}
    % Coined the term "Zero Trust"; challenged the concept of trusted networks.

    \subsubsection{Google's BeyondCorp}
    % One of the earliest large-scale real-world implementations.

\subsection{Evolution and Growing Industry Adoption}
% From niche concept to mainstream security framework.

\subsection{NIST SP 800-207 (2020)}
% Formal standardization of Zero Trust Architecture by NIST.

% ══════════════════════════════════════════════════════════════════════════════
% III. CORE PRINCIPLES OF ZERO TRUST  (~1.5 pages)
% ══════════════════════════════════════════════════════════════════════════════
\section{Core Principles of Zero Trust}

\subsection{``Never Trust, Always Verify''}

    \subsubsection{Continuous Authentication and Authorization}
    % Every request must be authenticated regardless of origin.

    \subsubsection{No Implicit Trust Based on Network Location}
    % Being inside the network no longer grants automatic access.

\subsection{Least-Privilege Access}

    \subsubsection{Users and Devices Get Only the Access They Need}
    % Role-based and attribute-based access control.

    \subsubsection{Minimizing the Blast Radius of a Breach}
    % How least privilege limits damage from compromised accounts.

\subsection{Assume Breach Mentality}

    \subsubsection{Designing Systems as if Attackers Are Already Inside}
    % Shifts focus from prevention-only to detection and containment.

    \subsubsection{Micro-Segmentation of Networks}
    % Isolating workloads to prevent lateral movement.

\subsection{Continuous Monitoring and Validation}

    \subsubsection{Real-Time Telemetry and Behavioral Analytics}
    % Logging, SIEM integration, anomaly detection.

    \subsubsection{Dynamic Policy Enforcement}
    % Policies adapt in real time based on risk signals.

% ══════════════════════════════════════════════════════════════════════════════
% IV. KEY COMPONENTS \& TECHNOLOGIES  (~1.5 pages)
% ══════════════════════════════════════════════════════════════════════════════
\section{Key Components and Technologies}

\subsection{Identity and Access Management (IAM)}

    \subsubsection{Multi-Factor Authentication (MFA)}
    % Why passwords alone are insufficient.

    \subsubsection{Single Sign-On (SSO) and Identity Providers}
    % Centralized identity management across services.

\subsection{Device Health and Endpoint Security}

    \subsubsection{Device Trust Scoring / Posture Assessment}
    % Evaluating device compliance before granting access.

    \subsubsection{Mobile Device Management (MDM)}
    % Enforcing policy on employee and BYOD devices.

\subsection{Micro-Segmentation}

    \subsubsection{Software-Defined Perimeters}
    % Creating dynamic, identity-bound network boundaries.

    \subsubsection{East-West Traffic Control}
    % Controlling lateral traffic inside the network.

\subsection{Policy Decision Point (PDP) and Policy Enforcement Point (PEP)}

    \subsubsection{How ZTA Makes Access Decisions}
    % The policy engine evaluates context and grants/denies access.

    \subsubsection{Role of the Policy Engine}
    % Coordinates trust signals from identity, device, and data sources.

\subsection{Encryption and Data Security}

    \subsubsection{Encryption in Transit and at Rest}
    % TLS, end-to-end encryption requirements.

    \subsubsection{Data Loss Prevention (DLP) Tools}
    % Preventing sensitive data from leaving controlled environments.

% ══════════════════════════════════════════════════════════════════════════════
% V. REAL-WORLD IMPLEMENTATIONS \& CASE STUDIES  (~1.5 pages)
% ══════════════════════════════════════════════════════════════════════════════
\section{Real-World Implementations and Case Studies}

\subsection{Google BeyondCorp}

    \subsubsection{Transition Away from VPN-Based Access}
    % Replaced traditional VPN with context-aware access.

    \subsubsection{Results and Lessons Learned}
    % Improved security posture; template for enterprise adoption.

\subsection{U.S. Federal Government Mandate}

    \subsubsection{Executive Order 14028 (2021)}
    % Required federal agencies to adopt ZTA frameworks.

    \subsubsection{CISA's Zero Trust Maturity Model}
    % Five pillars: Identity, Devices, Networks, Applications, Data.

\subsection{Enterprise Adoption Example}
% TODO: Choose a case study — e.g., Microsoft, Okta, a healthcare or
% financial services organization.

    \subsubsection{Implementation Overview}
    % What they deployed, timeline, tools used.

    \subsubsection{Outcomes}
    % Reduced breach surface, improved compliance posture.

\subsection{Common Implementation Challenges Observed Across Cases}
% Tie together patterns from all three examples above.

% ══════════════════════════════════════════════════════════════════════════════
% VI. CHALLENGES \& CRITICISMS  (~1 page)
% ══════════════════════════════════════════════════════════════════════════════
\section{Challenges and Criticisms}

\subsection{Implementation Complexity}

    \subsubsection{Legacy System Integration Issues}
    % Older infrastructure often cannot support ZTA natively.

    \subsubsection{High Upfront Cost and Resource Demands}
    % Budget and staffing requirements for full deployment.

\subsection{Organizational and Cultural Resistance}

    \subsubsection{User Friction from Stricter Authentication}
    % MFA fatigue; user experience trade-offs.

    \subsubsection{IT Staff Retraining Requirements}
    % Shift in mindset and skillset needed for ZTA operations.

\subsection{Vendor Lock-In Risks}

    \subsubsection{Reliance on Specific Platforms}
    % Risk of dependence on a single vendor ecosystem.

    \subsubsection{Interoperability Concerns}
    % Challenges integrating multi-vendor ZTA components.

\subsection{ZTA Is Not a Silver Bullet}

    \subsubsection{What It Does Not Protect Against}
    % e.g., insider threats with valid credentials, social engineering.

    \subsubsection{Ongoing Maintenance Requirements}
    % ZTA requires continuous tuning — it is not set-and-forget.

% ══════════════════════════════════════════════════════════════════════════════
% VII. FUTURE OUTLOOK  (~0.5 page)
% ══════════════════════════════════════════════════════════════════════════════
\section{Future Outlook}

\subsection{AI and ML Integration with ZTA}
% Adaptive, intelligent policy enforcement using behavioral models.

\subsection{ZTA in IoT and Operational Technology (OT) Environments}
% Extending Zero Trust beyond traditional IT to physical/industrial systems.

\subsection{Post-Quantum Cryptography Readiness}
% How ZTA frameworks must adapt to quantum computing threats.

\subsection{Regulatory Pressure Driving Broader Adoption}
% Increasing compliance requirements (CMMC, NIS2, etc.) mandating ZTA.

% ══════════════════════════════════════════════════════════════════════════════
% VIII. CONCLUSION  (~0.5 page)
% ══════════════════════════════════════════════════════════════════════════════
\section{Conclusion}

\subsection{Restate the Problem}
% Why the traditional perimeter model is no longer viable.

\subsection{Summarize ZTA's Core Value Proposition}
% Identity-centric, continuous verification, least privilege.

\subsection{Acknowledge Limitations}
% Honest assessment of cost, complexity, and gaps.

\subsection{Closing Argument}
% ZTA is not optional — it is the foundational security model for the
% modern enterprise.

% ══════════════════════════════════════════════════════════════════════════════
% REFERENCES
% ══════════════════════════════════════════════════════════════════════════════
\newpage
\bibliographystyle{apalike}   % Change to ieeetr, plain, etc. as needed
% \bibliography{references}   % Uncomment and point to your .bib file

% ── Suggested sources (add to your .bib file) ─────────────────────────────────
% Rose, S., Borchert, O., Mitchell, S., & Connelly, S. (2020).
%   Zero Trust Architecture. NIST Special Publication 800-207.
%   https://doi.org/10.6028/NIST.SP.800-207
%
% Kindervag, J. (2010). No More Chewy Centers: Introducing the Zero Trust
%   Model of Information Security. Forrester Research.
%
% Ward, B., & Beyer, B. (2014). BeyondCorp: A New Approach to Enterprise
%   Security. USENIX ;login: Magazine, 39(6).
%
% Executive Order 14028 on Improving the Nation's Cybersecurity. (2021).
%   Federal Register. https://www.federalregister.gov/d/2021-10460
%
% CISA. (2023). Zero Trust Maturity Model (Version 2.0).
%   Cybersecurity and Infrastructure Security Agency.
%   https://www.cisa.gov/zero-trust-maturity-model

\end{document}
